\documentclass[11pt,a4paper]{article}
\usepackage[utf8]{inputenc}
\usepackage[T1]{fontenc}
\usepackage{amsmath,amssymb}
\usepackage{geometry}
\geometry{margin=2.5cm}
\usepackage{graphicx}
\usepackage{hyperref}
\usepackage{listings}
\usepackage{xcolor}

\lstset{
  language=Python,
  basicstyle=\ttfamily\small,
  keywordstyle=\color{blue}\bfseries,
  commentstyle=\color{gray}\itshape,
  stringstyle=\color{orange},
  frame=single,
  breaklines=true,
  numbers=left,
  numberstyle=\tiny\color{gray}
}

\title{Manual Validation \& Debug Tooling\\
\large RTS Core Systems GUI Architecture}
\author{Senior UI/UX Engineering Team}
\date{\today}

\begin{document}
\maketitle
\tableofcontents
\newpage

\section{Introduction}

Validating the \texttt{ActionGateway} and \texttt{IUnitAction} queue directly from the core game loop is critical for ensuring non-functional requirement compliance. To achieve this without relying on complex JSON plans generated by the AI Team, we have designed a **Developer Debug Window**. This lightweight GUI acts as a manual override to the API, allowing developers to inject commands, simulate multi-player constraints, and observe asynchronous task completion via the \texttt{SenseAPI}.

\section{Architecture \& Decoupling}

The Developer Debug Window (\texttt{DebugCanvasLayer}) is constructed entirely programmatically within GDScript. This design choice prevents visual node corruption or UID misalignment within the editor environment.

\subsection{Design Principles}
\begin{enumerate}
    \item \textbf{Standalone Execution}: The \texttt{DebugWindow} interacts exclusively with the \texttt{ActionGateway} singleton. If the debug layer is omitted from production builds, the core game loop remains completely unaffected.
    \item \textbf{Live Telemetry Validation}: By querying the \texttt{SenseAPI} dynamically using \texttt{get\_unit\_status()}, developers can verify the internal transition from \texttt{IDLE} to \texttt{BUSY} in real-time.
    \item \textbf{Ownership Subversion Checks}: The UI contains an "Active Player ID" toggle. Injecting a command for Player 2's unit while spoofing Player 1's ID natively stress-tests the \texttt{[OWNERSHIP\_ERR]} rejection pathway.
\end{enumerate}

\section{Integration with ActionGateway}

\subsection{Command Dispatch}
User interactions are packaged and routed through the identical high-level wrappers exposed to the AI module.

\begin{lstlisting}[language=Python,caption={Example: Dispatching Composite Chains}]
func _on_btn_chop() -> void:
    var res = ActionGateway.go_chop_tree_and_return(
        unit_id, tree_id, active_player_id
    )
    if res:
        _log("[CMD_SENT] Command dispatched.")
\end{lstlisting}

\subsection{Telemetry Processing}
The sense feed verifies state transitions. For example, when a unit begins a composite command, the telemetry reads:
$$
\text{State: } \texttt{BUSY} \quad | \quad \text{Action: } \texttt{MOVE\_TO\_HARVEST}
$$
As the internal \texttt{queue\_empty} signal advances the state internally, the debug user sees the action shift seamlessly from \texttt{MOVE} to \texttt{HARVEST} to \texttt{RETURN}.

\section{Feedback Logging}

The tool completely negates the need for external system consoles by injecting event logs directly into a \texttt{RichTextLabel}. The following tags rigorously validate system integrity:

\begin{itemize}
    \item \texttt{[CMD\_SENT]}: Validation passed; command resides in the \texttt{CommandQueue}.
    \item \texttt{[CMD\_EXECUTING]}: The unit's \texttt{\_physics\_process} successfully popped and initiated the contract.
    \item \texttt{[CMD\_DENIED: WRONG\_OWNER]}: Successful verification that the Gateway actively blocks unauthorized manipulation.
    \item \texttt{[CMD\_COMPLETE]}: The \texttt{task\_completed} signal propagated correctly back out of the queue instance.
\end{itemize}

\section{Sandbox Analysis: Verification of Tree Integration}

To confirm compliance with the SEA-project architectural mandate (specifically the pathfinding integration identified for static environment objects), the Debug Window executes the following analytical steps when ``Chop Tree'' is targeted:

\begin{enumerate}
    \item \textbf{Dispatch Evaluation}: The `go\_chop` command enqueues a `UnitActionMove` targeted to the globally identified `tree\_id`.
    \item \textbf{Pathfind Acquisition}: The navigation server traces a path intersecting the outer bounds of the tree's \texttt{StaticMapObject} collision structure.
    \item \textbf{Fallback to Interaction}: Upon arrival, the gateway transitions cleanly, validating that target references are preserved correctly across multi-step sequences.
\end{enumerate}

\end{document}
